\documentclass[11pt]{article}

\usepackage[margin=1in]{geometry}
\usepackage{amsmath, amssymb}
\usepackage{booktabs}
\usepackage{array}
\usepackage{hyperref}
\usepackage{xcolor}
\usepackage{enumitem}
\usepackage{listings}

\hypersetup{
    colorlinks=true,
    linkcolor=blue,
    urlcolor=blue
}

\lstdefinestyle{terminal}{
    basicstyle=\ttfamily\small,
    breaklines=true,
    frame=single,
    columns=fullflexible,
    backgroundcolor=\color{gray!8}
}

\lstset{style=terminal}

\title{Reactor Engineering Engine: Capabilities, Inputs, and Worked Command Examples}
\author{Reactor Engine User Guide}
\date{\today}

\begin{document}
\maketitle

\tableofcontents
\newpage

\section{Purpose of the Code}

The main file is:

\begin{lstlisting}
reactor_engine.py
\end{lstlisting}

It is a Python command-line tool for common chemical reaction engineering calculations. The code can compute reactor design, conversion profiles, temperature profiles, reaction selectivity, reaction yield, kinetic parameters from experimental data, and parallel-reaction CSTR optimization.

The general command pattern is:

\begin{lstlisting}
python3 reactor_engine.py --mode MODE_NAME [inputs]
\end{lstlisting}

The most important input is always \texttt{--mode}, because it tells the code what type of problem to solve.

\section{Before Running}

From Terminal, clone the project and go to the project folder:

\begin{lstlisting}
git clone https://github.com/77bradbf/reactor-engine.git
cd reactor-engine
\end{lstlisting}

Install the base Python packages if needed:

\begin{lstlisting}
pip install -r requirements.txt
\end{lstlisting}

For advanced real-data/property calculations, install:

\begin{lstlisting}
pip install -r requirements-advanced.txt
\end{lstlisting}

Show all available options:

\begin{lstlisting}
python3 reactor_engine.py --help
\end{lstlisting}

Run without plots:

\begin{lstlisting}
python3 reactor_engine.py --mode series-batch --no-show
\end{lstlisting}

Run with plots:

\begin{lstlisting}
python3 reactor_engine.py --mode series-batch
\end{lstlisting}

Save plots instead of opening them:

\begin{lstlisting}
python3 reactor_engine.py --mode series-batch --save-plots plots --no-show
\end{lstlisting}

\section{Modes Available}

\begin{center}
\begin{tabular}{p{0.28\linewidth} p{0.62\linewidth}}
\toprule
\textbf{Mode} & \textbf{What It Computes} \\
\midrule
\texttt{cstr} & Single-reaction CSTR conversion or required volume. \\
\texttt{pfr} & Single-reaction isothermal PFR conversion profile. \\
\texttt{nonisothermal-pfr} & Coupled conversion and temperature profile for a PFR. \\
\texttt{series-batch} & Batch reactor for first-order series reactions $A \to B \to C$. \\
\texttt{series-cstr} & CSTR for first-order series reactions $A \to B \to C$. \\
\texttt{kinetic-fit-power-law} & Fits $r = kx_1^\alpha x_2^\beta$ from experimental data. \\
\texttt{parallel-cstr} & Optimizes selectivity and profit for parallel CSTR reactions. \\
\bottomrule
\end{tabular}
\end{center}

\section{Mode 1: Single-Reaction CSTR}

\subsection{What the Code Computes}

For a CSTR, the code uses the design equation:

\begin{equation}
V = \frac{F_{A0}X}{-r_A}
\end{equation}

The default reaction rate is second-order:

\begin{equation}
-r_A = k C_A C_B
\end{equation}

The default Arrhenius equation is:

\begin{equation}
k = A\exp\left(-\frac{E_a}{RT}\right)
\end{equation}

with:

\begin{align*}
A &= 5.7\times10^5 \\
E_a &= 33400\ \text{J/mol} \\
R &= 8.314\ \text{J/(mol K)}
\end{align*}

\subsection{Minimum Inputs}

For a CSTR, you must know either:

\begin{itemize}
    \item target conversion $X$, to calculate volume, or
    \item reactor volume $V$, to calculate conversion.
\end{itemize}

Useful inputs are:

\begin{itemize}
    \item \texttt{--target-x}
    \item \texttt{--volume}
    \item \texttt{--temperature}
    \item \texttt{--fa0}
    \item \texttt{--fb0}
    \item \texttt{--v0}
\end{itemize}

\subsection{Example 1A: Find CSTR Volume for a Target Conversion}

\textbf{Problem statement.} A liquid-phase reaction follows:

\begin{equation}
-r_A = k C_A C_B
\end{equation}

A feed contains $F_{A0}=10$ mol/min and $F_{B0}=15$ mol/min. The inlet volumetric flow rate is $v_0=100$ L/min. The reactor is isothermal at 360 K. Find the CSTR volume required for 65\% conversion.

\textbf{Command.}

\begin{lstlisting}
python3 reactor_engine.py --mode cstr --target-x 0.65 --temperature 360 --fa0 10 --fb0 15 --v0 100 --no-show
\end{lstlisting}

\textbf{What each input means.}

\begin{center}
\begin{tabular}{ll}
\toprule
\textbf{Input} & \textbf{Meaning} \\
\midrule
\texttt{--mode cstr} & Use CSTR design equation. \\
\texttt{--target-x 0.65} & Desired conversion is 65\%. \\
\texttt{--temperature 360} & Reactor temperature is 360 K. \\
\texttt{--fa0 10} & Inlet molar flow of A. \\
\texttt{--fb0 15} & Inlet molar flow of B. \\
\texttt{--v0 100} & Inlet volumetric flow rate. \\
\texttt{--no-show} & Print output without opening a plot window. \\
\bottomrule
\end{tabular}
\end{center}

\subsection{Example 1B: Find CSTR Conversion for a Known Volume}

\textbf{Problem statement.} The same reaction is run in a 250 L CSTR at 360 K. Find the exit conversion.

\textbf{Command.}

\begin{lstlisting}
python3 reactor_engine.py --mode cstr --volume 250 --temperature 360 --fa0 10 --fb0 15 --v0 100 --no-show
\end{lstlisting}

\section{Mode 2: Isothermal PFR}

\subsection{What the Code Computes}

For an isothermal plug flow reactor, the code integrates:

\begin{equation}
\frac{dX}{dV}=\frac{-r_A}{F_{A0}}
\end{equation}

The default rate law is:

\begin{equation}
-r_A = kC_A C_B
\end{equation}

\subsection{Minimum Inputs}

At minimum, use:

\begin{itemize}
    \item \texttt{--mode pfr}
    \item \texttt{--volume}
\end{itemize}

Add feed and temperature data when the problem gives them.

\subsection{Example 2A: Liquid-Phase PFR Conversion}

\textbf{Problem statement.} A second-order liquid-phase reaction is carried out in a 400 L PFR at 375 K. The feed contains $F_{A0}=8$ mol/min, $F_{B0}=12$ mol/min, and $v_0=80$ L/min. Find the exit conversion.

\textbf{Command.}

\begin{lstlisting}
python3 reactor_engine.py --mode pfr --volume 400 --temperature 375 --fa0 8 --fb0 12 --v0 80 --no-show
\end{lstlisting}

\subsection{Example 2B: PFR with a Different Temperature}

\textbf{Problem statement.} Repeat the PFR calculation at 390 K to see the effect of temperature on conversion.

\textbf{Command.}

\begin{lstlisting}
python3 reactor_engine.py --mode pfr --volume 400 --temperature 390 --fa0 8 --fb0 12 --v0 80 --no-show
\end{lstlisting}

\section{Mode 3: Gas-Phase PFR with Expansion}

\subsection{What the Code Computes}

For gas-phase reactions, volumetric flow may change with conversion. The code can use:

\begin{equation}
v = v_0(1+\epsilon X)
\end{equation}

where:

\begin{equation}
\epsilon = y_{A0}\delta
\end{equation}

\subsection{Minimum Inputs}

Use the normal PFR inputs plus:

\begin{itemize}
    \item \texttt{--gas-phase}
    \item \texttt{--y-a0}
    \item \texttt{--delta}
\end{itemize}

\subsection{Example 3A: Gas Expansion Reduces Concentration}

\textbf{Problem statement.} A gas-phase reaction is carried out in a 400 L PFR at 375 K. The inlet mole fraction of A is $y_{A0}=0.4$ and $\delta=1.5$. Find the conversion when expansion is included.

\textbf{Command.}

\begin{lstlisting}
python3 reactor_engine.py --mode pfr --volume 400 --temperature 375 --fa0 8 --fb0 12 --v0 80 --gas-phase --y-a0 0.4 --delta 1.5 --no-show
\end{lstlisting}

\subsection{Example 3B: Stronger Gas Expansion}

\textbf{Problem statement.} Repeat the gas-phase PFR problem with $\delta=2.5$.

\textbf{Command.}

\begin{lstlisting}
python3 reactor_engine.py --mode pfr --volume 400 --temperature 375 --fa0 8 --fb0 12 --v0 80 --gas-phase --y-a0 0.4 --delta 2.5 --no-show
\end{lstlisting}

\section{Mode 4: Non-Isothermal PFR}

\subsection{What the Code Computes}

The code solves the coupled PFR material and energy balance:

\begin{align}
\frac{dX}{dV} &= \frac{-r_A}{F_{A0}} \\
\frac{dT}{dV} &= \frac{Ua(T_a-T)+(-r_A)(-\Delta H_{rxn})}{\sum F_i C_{p,i}}
\end{align}

\subsection{Minimum Inputs}

Use:

\begin{itemize}
    \item \texttt{--mode nonisothermal-pfr}
    \item \texttt{--volume}
    \item \texttt{--t0}
    \item \texttt{--delta-h}
    \item \texttt{--ua}
    \item \texttt{--ta}
    \item heat capacities if different from defaults
\end{itemize}

\subsection{Example 4A: Exothermic PFR with Cooling}

\textbf{Problem statement.} An exothermic liquid-phase reaction is carried out in a 300 L PFR. Feed enters at 340 K with $F_{A0}=12$ mol/min, $F_{B0}=18$ mol/min, and $v_0=120$ L/min. The heat of reaction is $-45000$ J/mol, $Ua=350$, and coolant temperature is 320 K. Find conversion and exit temperature.

\textbf{Command.}

\begin{lstlisting}
python3 reactor_engine.py --mode nonisothermal-pfr --volume 300 --fa0 12 --fb0 18 --v0 120 --t0 340 --delta-h -45000 --ua 350 --ta 320 --cp-a 95 --cp-b 110 --cp-p 120 --no-show
\end{lstlisting}

\subsection{Example 4B: Weak Cooling}

\textbf{Problem statement.} Repeat the previous problem with weaker heat removal, $Ua=100$.

\textbf{Command.}

\begin{lstlisting}
python3 reactor_engine.py --mode nonisothermal-pfr --volume 300 --fa0 12 --fb0 18 --v0 120 --t0 340 --delta-h -45000 --ua 100 --ta 320 --cp-a 95 --cp-b 110 --cp-p 120 --no-show
\end{lstlisting}

\subsection{Example 4C: Gas-Phase Non-Isothermal PFR with Temperature Correction}

\textbf{Problem statement.} A gas-phase non-isothermal PFR has both conversion expansion and temperature expansion. Use $v=v_0(1+\epsilon X)(T/T_0)$.

\textbf{Command.}

\begin{lstlisting}
python3 reactor_engine.py --mode nonisothermal-pfr --volume 250 --fa0 10 --fb0 15 --v0 100 --t0 350 --gas-phase --y-a0 0.5 --delta 1.0 --ideal-gas-temperature-correction --delta-h -30000 --ua 250 --ta 330 --no-show
\end{lstlisting}

\section{Mode 5: Series Batch Reactor, $A \to B \to C$}

\subsection{What the Code Computes}

For first-order series reactions:

\begin{equation}
A \xrightarrow{k_1} B \xrightarrow{k_2} C
\end{equation}

The batch reactor expressions are:

\begin{align}
C_A(t) &= C_{A0}e^{-k_1t} \\
C_B(t) &= \frac{C_{A0}k_1}{k_2-k_1}\left(e^{-k_1t}-e^{-k_2t}\right) \\
C_C(t) &= C_{A0}-C_A(t)-C_B(t)
\end{align}

The time for maximum $B$ is:

\begin{equation}
t_{max} = \frac{\ln(k_1/k_2)}{k_1-k_2}
\end{equation}

\subsection{Minimum Inputs}

Use:

\begin{itemize}
    \item \texttt{--mode series-batch}
    \item \texttt{--series-k1}
    \item \texttt{--series-k2}
    \item \texttt{--ca0}
\end{itemize}

\subsection{Example 5A: Intermediate Product in Batch Reactor}

\textbf{Problem statement.} The elementary liquid-phase series reaction $A \to B \to C$ is carried out in a batch reactor. Let $k_1=0.5$ h$^{-1}$, $k_2=0.2$ h$^{-1}$, and $C_{A0}=2$ mol/L. Plot concentrations from 0 to 12 h and find the time for maximum $B$.

\textbf{Command.}

\begin{lstlisting}
python3 reactor_engine.py --mode series-batch --series-k1 0.5 --series-k2 0.2 --ca0 2 --time-start 0 --time-end 12
\end{lstlisting}

Use \texttt{--no-show} if you do not want the plot:

\begin{lstlisting}
python3 reactor_engine.py --mode series-batch --series-k1 0.5 --series-k2 0.2 --ca0 2 --time-start 0 --time-end 12 --no-show
\end{lstlisting}

\subsection{Example 5B: Faster Desired Step, Slower Undesired Step}

\textbf{Problem statement.} A batch reactor has $k_1=0.8$ h$^{-1}$ and $k_2=0.15$ h$^{-1}$. The initial A concentration is 3 mol/L. Find the best time to stop the batch to maximize $B$.

\textbf{Command.}

\begin{lstlisting}
python3 reactor_engine.py --mode series-batch --series-k1 0.8 --series-k2 0.15 --ca0 3 --time-start 0 --time-end 20 --no-show
\end{lstlisting}

\subsection{Example 5C: Undesired Step Is Fast}

\textbf{Problem statement.} A batch reactor has $k_1=0.4$ h$^{-1}$ and $k_2=1.1$ h$^{-1}$. Since $B$ disappears quickly, find the maximum possible concentration of $B$.

\textbf{Command.}

\begin{lstlisting}
python3 reactor_engine.py --mode series-batch --series-k1 0.4 --series-k2 1.1 --ca0 2.5 --time-start 0 --time-end 10 --no-show
\end{lstlisting}

\section{Mode 6: Series CSTR, $A \to B \to C$}

\subsection{What the Code Computes}

For a CSTR with first-order series reactions:

\begin{align}
C_A(\tau) &= \frac{C_{A0}}{1+k_1\tau} \\
C_B(\tau) &= \frac{k_1\tau C_A(\tau)}{1+k_2\tau} \\
C_C(\tau) &= C_{A0}-C_A(\tau)-C_B(\tau)
\end{align}

The residence time that maximizes $B$ is:

\begin{equation}
\tau_{opt}=\frac{1}{\sqrt{k_1k_2}}
\end{equation}

\subsection{Example 6A: CSTR Intermediate Optimization}

\textbf{Problem statement.} The same reaction $A \to B \to C$ is carried out in a CSTR. Use $k_1=0.5$ h$^{-1}$, $k_2=0.2$ h$^{-1}$, and $C_{A0}=2$ mol/L. Find the residence time that maximizes $B$.

\textbf{Command.}

\begin{lstlisting}
python3 reactor_engine.py --mode series-cstr --series-k1 0.5 --series-k2 0.2 --ca0 2 --time-start 0 --time-end 12 --no-show
\end{lstlisting}

\subsection{Example 6B: Larger Initial Concentration}

\textbf{Problem statement.} A CSTR uses the same rate constants but $C_{A0}=5$ mol/L. Find the maximum $B$ concentration and the optimum residence time.

\textbf{Command.}

\begin{lstlisting}
python3 reactor_engine.py --mode series-cstr --series-k1 0.5 --series-k2 0.2 --ca0 5 --time-start 0 --time-end 12 --no-show
\end{lstlisting}

\subsection{Example 6C: Different Rate Constants}

\textbf{Problem statement.} A CSTR has $k_1=1.2$ h$^{-1}$ and $k_2=0.3$ h$^{-1}$. The feed concentration of A is 1.8 mol/L. Find the residence time that maximizes $B$.

\textbf{Command.}

\begin{lstlisting}
python3 reactor_engine.py --mode series-cstr --series-k1 1.2 --series-k2 0.3 --ca0 1.8 --time-start 0 --time-end 10 --no-show
\end{lstlisting}

\section{Mode 7: Power-Law Kinetic Fitting}

\subsection{What the Code Computes}

This mode fits a two-reactant power-law rate model:

\begin{equation}
r = kx_1^\alpha x_2^\beta
\end{equation}

Taking logs gives:

\begin{equation}
\ln r = \ln k + \alpha \ln x_1 + \beta \ln x_2
\end{equation}

The code uses linear least squares on the log-transformed data to calculate:

\begin{itemize}
    \item $\alpha$
    \item $\beta$
    \item $k$
    \item $R^2$ on $\ln(r)$
    \item predicted rates
    \item an observed-versus-predicted plot
\end{itemize}

\subsection{Minimum Inputs}

If using the built-in photo example, only this is needed:

\begin{lstlisting}
python3 reactor_engine.py --mode kinetic-fit-power-law --no-show
\end{lstlisting}

For new data, use:

\begin{itemize}
    \item \texttt{--fit-data-csv examples/myfile.csv}, or
    \item \texttt{--fit-data-csv examples/myfile.xlsx}
\end{itemize}

The option is called \texttt{--fit-data-csv}, but it accepts both CSV and Excel files.

\subsection{Data Format A: Rate Already Given}

Use this table if the problem directly gives the measured rate:

\begin{center}
\begin{tabular}{ccc}
\toprule
\texttt{x1} & \texttt{x2} & \texttt{rate} \\
\midrule
0.20 & 0.30 & 0.0048 \\
0.40 & 0.30 & 0.0096 \\
0.20 & 0.60 & 0.0068 \\
0.50 & 0.50 & 0.0141 \\
0.80 & 0.40 & 0.0202 \\
\bottomrule
\end{tabular}
\end{center}

Save this as:

\begin{lstlisting}
examples/power_fit_example.xlsx
\end{lstlisting}

or:

\begin{lstlisting}
examples/power_fit_example.csv
\end{lstlisting}

Then run:

\begin{lstlisting}
python3 reactor_engine.py --mode kinetic-fit-power-law --fit-data-csv examples/power_fit_example.xlsx --no-show
\end{lstlisting}

\subsection{Example 7A: Fit a New Catalytic Rate Law}

\textbf{Problem statement.} A catalytic reaction $A+B\to P$ follows:

\begin{equation}
r_P = kP_A^\alpha P_B^\beta
\end{equation}

The experimental data are shown in the previous table. Determine $\alpha$, $\beta$, and $k$.

\textbf{Command.}

\begin{lstlisting}
python3 reactor_engine.py --mode kinetic-fit-power-law --fit-data-csv examples/power_fit_example.xlsx --no-show
\end{lstlisting}

\subsection{Data Format B: Product Mole Fraction Given}

Use this table if the problem gives total outlet flow and product mole fraction:

\begin{center}
\begin{tabular}{cccc}
\toprule
\texttt{total\_flow} & \texttt{x1} & \texttt{x2} & \texttt{y\_product} \\
\midrule
0.61 & 0.50 & 0.51 & 0.07 \\
0.60 & 1.00 & 0.54 & 0.16 \\
0.30 & 0.40 & 0.60 & 0.16 \\
0.72 & 0.60 & 0.60 & 0.10 \\
0.52 & 0.60 & 0.40 & 0.06 \\
\bottomrule
\end{tabular}
\end{center}

The code calculates:

\begin{equation}
r = \frac{F_T y_P}{W_{cat}}
\end{equation}

where $F_T$ is total molar flow, $y_P$ is product mole fraction, and $W_{cat}$ is catalyst mass.

\subsection{Example 7B: Gas-Phase Catalytic CSTR Fit}

\textbf{Problem statement.} The irreversible gas-phase reaction

\begin{equation}
H_2 + C_2H_6 \to 2CH_4
\end{equation}

is carried out over nickel catalyst. The model is:

\begin{equation}
r'_{CH_4}=kP_{C_2H_6}^{\alpha}P_{H_2}^{\beta}
\end{equation}

A CSTR contains four baskets with 10 g catalyst each, so $W_{cat}=40$ g. Determine $\alpha$, $\beta$, and $k$.

\textbf{Command.}

\begin{lstlisting}
python3 reactor_engine.py --mode kinetic-fit-power-law --fit-data-csv examples/kinetic_fit_photo.csv --catalyst-weight 40 --no-show
\end{lstlisting}

\subsection{Example 7C: Same Fit with a Plot}

Remove \texttt{--no-show}:

\begin{lstlisting}
python3 reactor_engine.py --mode kinetic-fit-power-law --fit-data-csv examples/kinetic_fit_photo.csv --catalyst-weight 40
\end{lstlisting}

\subsection{Example 7D: Saving the Fit Plot}

\begin{lstlisting}
python3 reactor_engine.py --mode kinetic-fit-power-law --fit-data-csv examples/kinetic_fit_photo.csv --catalyst-weight 40 --save-plots plots --no-show
\end{lstlisting}

\section{Mode 8: Parallel-Reaction CSTR Optimization}

\subsection{What the Code Computes}

This mode handles parallel reactions:

\begin{align}
A+B &\to D &&\text{desired product} \\
A+B &\to U_1 &&\text{undesired product} \\
A+B &\to U_2 &&\text{undesired product}
\end{align}

The general rate laws are:

\begin{align}
r_D &= k_D C_A^{a_D}C_B^{b_D} \\
r_{U1} &= k_{U1} C_A^{a_{U1}}C_B^{b_{U1}} \\
r_{U2} &= k_{U2} C_A^{a_{U2}}C_B^{b_{U2}}
\end{align}

The instantaneous selectivity is:

\begin{equation}
S_{D/(U1+U2)} = \frac{r_D}{r_{U1}+r_{U2}}
\end{equation}

The code finds the exit $C_A$ that maximizes selectivity. It also estimates residence time, volumetric flow rate, and profit difference using old and optimized outlet concentrations.

\subsection{Minimum Inputs}

Important inputs are:

\begin{itemize}
    \item feed concentrations: \texttt{--parallel-ca0}, \texttt{--parallel-cb0}
    \item known exit B concentration: \texttt{--parallel-exit-cb}
    \item reactor volume: \texttt{--parallel-volume}
    \item rate constants: \texttt{--parallel-kd}, \texttt{--parallel-ku1}, \texttt{--parallel-ku2}
    \item reaction orders: \texttt{--d-exp-a}, \texttt{--d-exp-b}, etc.
    \item prices and costs if profit is required
\end{itemize}

\subsection{Example 8A: Parallel CSTR with Selectivity and Profit}

\textbf{Problem statement.} A 500 L CSTR is used for the liquid-phase parallel reactions:

\begin{align*}
A+B &\to D \\
A+B &\to U_1 \\
A+B &\to U_2
\end{align*}

The feed has $C_{A0}=2.5$ mol/L and $C_{B0}=1.8$ mol/L. The outlet concentration of B is 0.6 mol/L. The rate laws are:

\begin{align*}
r_D &= 0.45C_A^1C_B^1 \\
r_{U1} &= 0.12C_A^{0.5}C_B^1 \\
r_{U2} &= 0.08C_A^2C_B^1
\end{align*}

Product D sells for 75 dollars/mol, and disposal of each undesired product costs 15 dollars/mol. The previous outlet stream had $D=1.1$, $U_1=0.4$, $U_2=0.5$ mol/L. The optimized outlet stream has $D=1.4$, $U_1=0.25$, $U_2=0.35$ mol/L. Find the exit $C_A$ that maximizes selectivity, the residence time, volumetric flow rate, and profit increase.

\textbf{Command.}

\begin{lstlisting}
python3 reactor_engine.py --mode parallel-cstr --parallel-ca0 2.5 --parallel-cb0 1.8 --parallel-exit-cb 0.6 --parallel-volume 500 --parallel-kd 0.45 --parallel-ku1 0.12 --parallel-ku2 0.08 --d-exp-a 1 --d-exp-b 1 --u1-exp-a 0.5 --u1-exp-b 1 --u2-exp-a 2 --u2-exp-b 1 --desired-price 75 --undesired-cost 15 --previous-d 1.1 --previous-u1 0.4 --previous-u2 0.5 --optimized-d 1.4 --optimized-u1 0.25 --optimized-u2 0.35 --no-show
\end{lstlisting}

\subsection{Example 8B: Same Problem with a Plot}

Remove \texttt{--no-show}:

\begin{lstlisting}
python3 reactor_engine.py --mode parallel-cstr --parallel-ca0 2.5 --parallel-cb0 1.8 --parallel-exit-cb 0.6 --parallel-volume 500 --parallel-kd 0.45 --parallel-ku1 0.12 --parallel-ku2 0.08 --d-exp-a 1 --d-exp-b 1 --u1-exp-a 0.5 --u1-exp-b 1 --u2-exp-a 2 --u2-exp-b 1 --desired-price 75 --undesired-cost 15 --previous-d 1.1 --previous-u1 0.4 --previous-u2 0.5 --optimized-d 1.4 --optimized-u1 0.25 --optimized-u2 0.35
\end{lstlisting}

\subsection{Example 8C: Change the Undesired Reaction Order}

\textbf{Problem statement.} Suppose the second undesired reaction is less sensitive to A:

\begin{equation}
r_{U2}=0.08C_A^{1.5}C_B
\end{equation}

Find the new optimum.

\textbf{Command.}

\begin{lstlisting}
python3 reactor_engine.py --mode parallel-cstr --parallel-ca0 2.5 --parallel-cb0 1.8 --parallel-exit-cb 0.6 --parallel-volume 500 --parallel-kd 0.45 --parallel-ku1 0.12 --parallel-ku2 0.08 --d-exp-a 1 --d-exp-b 1 --u1-exp-a 0.5 --u1-exp-b 1 --u2-exp-a 1.5 --u2-exp-b 1 --no-show
\end{lstlisting}

\subsection{Example 8D: Higher Desired Product Price}

\textbf{Problem statement.} If product D sells for 120 dollars/mol instead of 75 dollars/mol, calculate the new profit increase.

\textbf{Command.}

\begin{lstlisting}
python3 reactor_engine.py --mode parallel-cstr --parallel-ca0 2.5 --parallel-cb0 1.8 --parallel-exit-cb 0.6 --parallel-volume 500 --parallel-kd 0.45 --parallel-ku1 0.12 --parallel-ku2 0.08 --d-exp-a 1 --d-exp-b 1 --u1-exp-a 0.5 --u1-exp-b 1 --u2-exp-a 2 --u2-exp-b 1 --desired-price 120 --undesired-cost 15 --previous-d 1.1 --previous-u1 0.4 --previous-u2 0.5 --optimized-d 1.4 --optimized-u1 0.25 --optimized-u2 0.35 --no-show
\end{lstlisting}

\section{Common Input Checklist}

\subsection{For Reactor Design Problems}

Ask:

\begin{enumerate}
    \item What reactor type is it: CSTR, PFR, batch, or non-isothermal PFR?
    \item What is the rate law?
    \item What are the rate constants and temperature?
    \item What are the feed molar flows or concentrations?
    \item Is the problem asking for conversion, volume, residence time, selectivity, yield, or temperature?
    \item Is it gas phase? If yes, is expansion important?
\end{enumerate}

\subsection{For Kinetic Fitting Problems}

Ask:

\begin{enumerate}
    \item Does the problem give measured rates directly?
    \item Or does it give flow rate and product mole fraction?
    \item What are $x_1$ and $x_2$? Usually pressures or concentrations.
    \item What is the catalyst weight, if rates must be calculated from product flow?
    \item Are all values positive? Log fitting requires positive data.
\end{enumerate}

\subsection{For Parallel Selectivity Problems}

Ask:

\begin{enumerate}
    \item Which product is desired?
    \item Which products are undesired?
    \item What are the rate laws for each path?
    \item Which concentration is fixed at the exit?
    \item What concentration should be optimized?
    \item Is profit required? If yes, what are product values and disposal costs?
\end{enumerate}

\section{All Major Command-Line Inputs}

\begin{center}
\begin{tabular}{p{0.34\linewidth} p{0.56\linewidth}}
\toprule
\textbf{Input} & \textbf{Meaning} \\
\midrule
\texttt{--mode} & Selects the calculation type. \\
\texttt{--volume} & Reactor volume. \\
\texttt{--target-x} & Target conversion for CSTR volume calculation. \\
\texttt{--temperature} & Isothermal reactor temperature. \\
\texttt{--fa0}, \texttt{--fb0} & Inlet molar flows of A and B. \\
\texttt{--v0} & Inlet volumetric flow rate. \\
\texttt{--arrhenius-a}, \texttt{--ea} & Arrhenius parameters. \\
\texttt{--gas-phase} & Turns on gas expansion. \\
\texttt{--y-a0}, \texttt{--delta} & Gas expansion parameters. \\
\texttt{--delta-h}, \texttt{--ua}, \texttt{--ta} & Energy balance parameters. \\
\texttt{--series-k1}, \texttt{--series-k2} & Series reaction rate constants. \\
\texttt{--ca0} & Initial or feed concentration of A for series reactions. \\
\texttt{--fit-data-csv} & CSV or Excel data file for kinetic fitting. \\
\texttt{--catalyst-weight} & Catalyst mass for product-flow data. \\
\texttt{--parallel-ca0}, \texttt{--parallel-cb0} & Feed concentrations for parallel CSTR. \\
\texttt{--parallel-kd}, \texttt{--parallel-ku1}, \texttt{--parallel-ku2} & Parallel reaction rate constants. \\
\texttt{--d-exp-a}, \texttt{--u1-exp-a}, etc. & Reaction orders in the parallel rate laws. \\
\texttt{--desired-price}, \texttt{--undesired-cost} & Economic values for profit calculation. \\
\texttt{--no-show} & Suppresses plot windows. \\
\texttt{--save-plots} & Saves plots to a folder. \\
\bottomrule
\end{tabular}
\end{center}

\section{Summary}

The code is most useful when you first classify the problem:

\begin{itemize}
    \item Use \texttt{cstr} for a single CSTR design equation.
    \item Use \texttt{pfr} for a single isothermal PFR profile.
    \item Use \texttt{nonisothermal-pfr} when temperature changes with volume.
    \item Use \texttt{series-batch} for batch $A\to B\to C$ problems.
    \item Use \texttt{series-cstr} for CSTR $A\to B\to C$ problems.
    \item Use \texttt{kinetic-fit-power-law} when the problem gives experimental data and asks for reaction orders and $k$.
    \item Use \texttt{parallel-cstr} when the problem asks for selectivity, yield, waste products, or profit optimization in parallel reactions.
\end{itemize}

\end{document}
